\documentclass[10pt,a4paper]{amsart}
\usepackage[margin=19mm]{geometry}
\usepackage{amsmath,amssymb,mathtools}
\usepackage[scr=rsfs,cal=euler]{mathalfa}
\usepackage{xfrac}
\usepackage[unicode,hidelinks,bookmarks=false]{hyperref}
\newcommand{\ie}{i.e.\ }
\newcommand{\cf}{cf.\ }
\newcommand{\resp}{respectively\ }
\newcommand{\Subsection}{\subsection}
\providecommand{\nobreakdash}{\nobreak\hbox{-}}
\setlength{\emergencystretch}{2em}
\multlinegap0pt
\usepackage{bm}
\usepackage{bbm}
\usepackage{dsfont}
\usepackage{xspace}
\usepackage{cancel}
\usepackage{mathtools}
\usepackage{amsthm}
\makeatletter
\@ifundefined{theorem}{\newtheorem{theorem}{Theorem}}{}
\@ifundefined{corollary}{\newtheorem{corollary}[theorem]{Corollary}}{}
\@ifundefined{lemma}{\newtheorem{lemma}[theorem]{Lemma}}{}
\@ifundefined{proposition}{\newtheorem{proposition}[theorem]{Proposition}}{}
\makeatother
\makeatletter
\expandafter\ifx\csname proof\endcsname\relax
\newenvironment{proof}[1][Proof]{\noindent\textit{#1.} }{}
\fi
\makeatother
\title[Affine Hulls and Simplices: a Constructive Analysis]{Affine Hulls and Simplices: a Constructive Analysis}
\author{Douglas S. Bridges}
\date{Source: arXiv:2509.20633v1 (CC BY 4.0)}
\begin{document}
\maketitle

\noindent\emph{Source and licence.} Affine Hulls and Simplices: a Constructive Analysis. Version arXiv:2509.20633v1 (CC BY 4.0), distributed under the Creative Commons Attribution 4.0 International licence, as recorded in the arXiv licence metadata for this exact version and shown on the abstract page https://arxiv.org/abs/2509.20633v1. Full rights notice: licenses/arxiv-2509.20633-CC-BY-4.0.txt.\par\smallskip

\noindent\textit{Editorial scope.} Counted unit: the Theorem whose source label is \texttt{sept13t1} in the article (source0.tex lines 1368--1376, source label \texttt{sept13t1}), delivered together with the article's own complete original proof of it (source0.tex lines 1378--1469, one \texttt{proof} environment of 92 lines). No mathematics is written, repaired or re-derived: statement and proof are the article's own text line for line. Required context retained uncounted: aug04l1A (lines 413--432); SS1 (lines 420--424); jul21p3 (lines 1124--1138); sept04c1 (lines 1297--1303); sept12l1 (lines 1345--1350). Every definition, display and label of those lines is retained unchanged. Omitted from this package and not redistributed: 1--412, 425--1123, 1139--1296, 1304--1344, 1351--1367, 1377--1377, 1470--1561. Nothing inside a retained excerpt is omitted or altered: every retained body is delivered exactly as the article's lines are written, so no omission receipt of any kind accompanies this package. Citations: the retained excerpts carry no citation command at all, so no bibliography entry of the article is transcribed and no reference is redistributed. Reference adaptation: none was needed and none was made; every reference inside the delivered unit resolves inside it, so no unresolved reference or citation is produced. Printed slips kept literal: no printed word, symbol, hypothesis or number of the counted statement or of its proof was altered. Layout and class: the article's own class is \texttt{article} with packages including amssymb, amsfonts, amsmath, euler, graphicx; the wrapper is the lane's pinned \texttt{amsart} editorial wrapper at 10pt on A4, which restates the theorem-like environments the retained text needs and adds the article's own macro definitions that the retained text needs (none), with extra wrapper packages (bm, bbm, dsfont, xspace, cancel, mathtools). The delivered numbering is the unit's own. Assets: no retained span contains a figure environment, a graphics-inclusion command, a PDF-page inclusion, a table or a TikZ picture; no image file is copied into this package, the package declares no asset dependency and no image file is redistributed with it. Licence: version arXiv:2509.20633v1 of this article is distributed under the Creative Commons Attribution 4.0 International licence, as recorded in the arXiv licence metadata for this exact version and shown on the abstract page https://arxiv.org/abs/2509.20633v1. A full rights notice is delivered at licenses/arxiv-2509.20633-CC-BY-4.0.txt. Characters: every retained excerpt is pure ASCII, so the delivered engine, XeLaTeX with its default Latin Modern text font, reports no missing character.
\begin{lemma}
\label{aug04l1A}Let $a_{1},\ldots,a_{n+1}$ be affinely independent elements of
a real normed linear space $X$. Then for each $\varepsilon>0$ there exists
$\delta>0$ with the following properties.

\begin{enumerate}
\item[\emph{(i)}] If%
\begin{equation}
\ \mathbf{\xi},\mathbf{\eta}\in\mathbf{R}^{n+1},\ \sum_{k=1}^{n+1}\xi
_{k}=1=\sum_{k=1}^{n+1}\eta_{k},\ x=\sum_{k=1}^{n+1}\xi_{k}a_{k}%
,\ y=\sum_{k=1}^{n+1}\eta_{k}a_{k}, \label{SS1}%
\end{equation}
$\ $and$\ \left\Vert x-y\right\Vert <\delta,$ then $\sum_{k=1}^{n+1}\left\vert
\xi_{k}-\eta_{k}\right\vert <\varepsilon$.

\item[\emph{(ii)}] If \emph{(\ref{SS1})} holds and $\sum_{k=1}^{n}\left\vert
\xi_{k}-\eta_{k}\right\vert <\delta$, then $\left\Vert x-y\right\Vert
<\varepsilon$.
\end{enumerate}
\end{lemma}

\begin{proposition}
\label{jul21p3}Let $\Sigma$ be an $n$-simplex with vertices $a_{1}%
,\ldots,a_{n+1}$ in a real normed linear space $X$, let $\mathbf{\lambda}%
\in\Delta_{n}$, and let $c=\sum_{k=1}^{n}\lambda_{k}a_{k}\in\Sigma$. Then the
following conditions are equivalent

\begin{enumerate}
\item[\emph{(i)}] $c$ is a relative interior point of $\Sigma$ in $X.$

\item[\emph{(ii)}] $c$ is bounded away from each $(n-1)$-face of $\Sigma$.

\item[\emph{(iii)}] $c=\sum_{k=1}^{n+1}\lambda_{k}a_{k}$, where
$\mathbf{\lambda}\in\Delta_{n}^{+}$.
\end{enumerate}
\end{proposition}

\begin{corollary}
\label{sept04c1}Let $a_{1},\ldots,a_{n+1}$ be affinely independent elements of
a real normed linear space $X$. Then there exists $\delta>0$ such that if
$x_{1},\ldots,x_{n+1}\ $belong to $X$ and $\left\Vert a_{k}-x_{k}\right\Vert
<\delta$ for each $k\leq n+1$, then the vectors $x_{1},\ldots,x_{n+1}$ are
affinely independent.
\end{corollary}

\begin{lemma}
\label{sept12l1}For each $\varepsilon>0$ there exists $\delta>0$ such that if
$A\equiv\left[  \xi_{kj}\right]  $ is an $n$-by-$n$ matrix over $\mathbf{R}$
and $\left\Vert A-I\right\Vert _{1}<\delta$, then $\det A>1/2\,\ $and
$\left\Vert A^{-1}-I\right\Vert _{1}<\varepsilon$.
\end{lemma}

\begin{theorem}
\label{sept13t1}Let $X$ be a real normed linear space, $\Sigma$ an $n$-simplex
with vertices $a_{1},\ldots,a_{n+1}$ in $X$, and $c\ $a relative interior
point of $\Sigma$. There exists $\delta>0$ such that if $x_{k}\in
\mathsf{aff}(\Sigma)$ and $\left\Vert x_{k}-a_{k}\right\Vert <\delta$ for each
$k\leq n+1$, then the points $x_{1},\ldots,x_{n+1}$ are affinely independent,
and$\ c$ is a relative interior point of the $n$-simplex $\mathsf{co}\left\{
x_{1},\ldots,x_{n+1}\right\}  $.
\end{theorem}

\begin{proof}
By Corollary \ref{sept04c1}, there exists $\delta_{1}$ with $0<\delta_{1}<1/2$
such that if $x_{1},\ldots,x_{n+1}\ $belong to $X$ and $\left\Vert a_{k}%
-x_{k}\right\Vert <\delta_{1}$ for each $k\leq n+1$, then $x_{1}%
,\ldots,x_{n+1}$ are affinely independent. By Proposition \ref{jul21p3},
$c=\sum_{k=1}^{n+1}\lambda_{k}a_{k}$, where $\mathbf{\lambda}\in\Delta_{n}%
^{+}$. Let%
\begin{equation}
m=\min_{1\leq k\leq n+1}\lambda_{k}.\nonumber
\end{equation}
With $\kappa$ as in the sentence immediately before Lemma \ref{sept12l1},
choose $\delta_{2}$ with $0<\delta_{2}<\delta_{1}$ such that if $A$ is an
$n$-by-$n$ matrix over $\mathbf{R}$ and $\left\Vert A-I\right\Vert _{1}%
<\delta_{2}$, then $\det A>1/2$, $\left\Vert A^{-1}-I\right\Vert
_{1}<m/2\kappa$, and therefore $\left\Vert A^{-1}-I\right\Vert _{0}<m/2$.
Using Lemma \ref{aug04l1A}, construct $\delta$ with $0<\delta<\delta_{2}$ such
that if%
\[
\ \mathbf{\xi},\mathbf{\eta}\in\mathbf{R}^{n+1},\ \sum_{k=1}^{n+1}\xi
_{k}=1=\sum_{k=1}^{n+1}\eta_{k},\ x=\sum_{k=1}^{n+1}\xi_{k}a_{k}%
,\ y=\sum_{k=1}^{n+1}\eta_{k}a_{k},
\]
and $\left\Vert \sum_{k=1}^{n+1}\left(  \xi_{k}-\eta_{k}\right)
a_{k}\right\Vert $ then $\left\vert \xi_{k}-\eta_{k}\right\vert <\delta_{2}$
for each $k\leq n+1$. Consider points $x_{1},\ldots,x_{n+1}$ of $\mathsf{aff}%
(\Sigma)$ such that $\left\Vert x_{k}-a_{k}\right\Vert <\delta$ for each
$k\leq n+1$. Since $\delta<\delta_{1}$, those points $x_{k}$ are affinely
independent, so $\Sigma^{\prime}\equiv\mathsf{co}\left\{  x_{1},\ldots
,x_{n+1}\right\}  $ is an $n$-simplex. For each $k$ there exists a unique
$\mathbf{\xi}_{k}\equiv\left(  \xi_{k,1},\xi_{k,2},\ldots,\xi_{k,n+1}\right)
\in\mathbf{R}^{n+1}$ such that $\sum_{j=1}^{n+1}\xi_{k,j}=1$ and $x_{k}%
=\sum_{j=1}^{n+1}\xi_{k,j}a_{j}$. Then%
\begin{equation}
\left\vert \xi_{k,k}-1\right\vert <\delta_{2}\text{ and }\left\vert \xi
_{k,j}\right\vert <\delta_{2}\text{ for }j\neq k. \label{B0}%
\end{equation}
I claim that the equations%
\begin{equation}
\sum_{k=1}^{n+1}\gamma_{k}\xi_{k,j}=\lambda_{j}\ \ \ \ (1\leq j\leq n+1).
\label{B1}%
\end{equation}
\label{000001}have a unique solution $\mathbf{\gamma}$, that $\mathbf{\gamma
}\in\Delta_{n}^{+}$, and that $c=\sum_{j=1}^{n+1}\gamma_{j}x_{j}$. To justify
this claim, letting $\Xi$ be the $n$-by-$n$ matrix with $\left(  j,k\right)
^{\mathsf{th}}$ component $\xi_{j,k}$, we see that the equations can be
written as $\Xi^{T}\mathbf{\gamma}=\mathbf{\lambda}$, where $^{T}$ denotes
transpose. From (\ref{B0}) we see that $\left\Vert \Xi^{T}-I\right\Vert
_{1}<\delta_{2}$, so
\begin{equation}
\det\Xi^{T}>\frac{1}{2}\text{ and }\left\Vert \left(  \Xi^{T}\right)
^{-1}-I\right\Vert _{0}<\frac{m}{2}. \label{B-1}%
\end{equation}
The equations (\ref{B1}) therefore have a unique solution $\mathbf{\gamma
}=\left(  \Xi^{T}\right)  ^{-1}\mathbf{\lambda}$. Then%
\begin{align*}
c  &  =\sum_{j=1}^{n+1}\lambda_{j}a_{j}=\sum_{j=1}^{n+1}\sum_{k=1}^{n+1}%
\gamma_{k}\xi_{k,j}a_{j}\\
&  =\sum_{k=1}^{n+1}\sum_{j=1}^{n+1}\gamma_{k}\xi_{k,j}a_{j}=\sum_{k=1}%
^{n+1}\gamma_{k}\sum_{j=1}^{n+1}\xi_{k,j}a_{j}=\sum_{k=1}^{n+1}\gamma_{k}x_{k}%
\end{align*}
and%
\begin{align*}
\sum_{k=1}^{n+1}\gamma_{k}  &  =\sum_{k=1}^{n+1}\gamma_{k}1=\sum_{k=1}%
^{n+1}\gamma_{k}\sum_{j=1}^{n+1}\xi_{k,j}=\sum_{k=1}^{n+1}\sum_{j=1}%
^{n+1}\gamma_{k}\xi_{k,j}\\
&  =\sum_{j=1}^{n+1}\sum_{k=1}^{n+1}\gamma_{k}\xi_{k,j}=\sum_{j=1}%
^{n+1}\lambda_{j}=1\text{,}%
\end{align*}
from which we see that $c\in\mathsf{aff}\left\{  x_{1},\ldots,x_{n+1}\right\}
$. On the other hand, by (\ref{B-1}),
\begin{align*}
\sum_{j=1}^{n+1}\left\vert \gamma_{j}-\lambda_{j}\right\vert  &  =\left\Vert
\left(  \Xi^{T}\right)  ^{-1}\mathbf{\lambda-\lambda}\right\Vert
_{1}=\left\Vert \left(  \left(  \Xi^{T}\right)  ^{-1}-I\right)
\mathbf{\lambda}\right\Vert _{1}\\
&  \leq\left\Vert \left(  \Xi^{T}\right)  ^{-1}-I\right\Vert _{0}\left\Vert
\mathbf{\lambda}\right\Vert _{1}=\left\Vert \left(  \Xi^{T}\right)
^{-1}-I\right\Vert _{0}<\frac{m}{2},
\end{align*}
so for each $j\leq n+1$,%
\[
\gamma_{j}>\lambda_{j}-\frac{m}{2}\geq m-\frac{m}{2}>0\text{.}%
\]
Since also $\sum_{j=1}^{n+1}\gamma_{j}=1$, $\mathbf{\gamma}\in\Delta_{n}^{+}$
and therefore, by Proposition \ref{jul21p3}(iii), $c$ is a relative interior
point $\Sigma^{\prime}$.%
%TCIMACRO{\TeXButton{hfill}{\hfill}}%
%BeginExpansion
\hfill
%EndExpansion
$\square$
\end{proof}

\end{document}
